\section{Stationary source spectroscopy beyond finite-rank Floquet selection}
\label{sec:algebraic-stationary-completion}

The Haar obstruction and the finite-dimensional Riesz construction refer to
different spectral domains. They can be reconciled at the level of stationary
source correlations without identifying their state prescriptions. The
construction below uses standard invariant-state and GNS methods. Its role
here is to extend the source observable, not to establish a new general GNS
theorem or to assume a physical vacuum. The microscopic dynamics, the charged
source algebra, and the selected limiting functional remain physical inputs.

\subsection{The precise obstruction}

The commuting-character argument in the Haar analysis excludes nonzero
finite-rank Riesz projections of the microscopic update under its stated
commutation and invariant-subspace hypotheses. It also excludes positive
trace-class operators commuting with the character translation. It does not
exclude every density commuting with the update, because commutation with the
update does not imply commutation with that translation. In particular, the
Haar argument must not be restated as a theorem that the full update has purely
continuous spectrum. Infinite-multiplicity point spectrum remains logically
possible.

A separate useful fact is that a unitary with no eigenvectors admits no
nonzero positive trace-class operator commuting with it. Indeed, a nonzero
positive compact operator has a positive eigenvalue with a finite-dimensional
eigenspace. Commutation makes that eigenspace invariant under the unitary,
which would then have an eigenvector. This fact explains why a stationary
state must sometimes be formulated as a functional on an observable algebra
rather than as a trace-class density on the original Hilbert space.

\subsection{Stationary functionals and their selection}

Let $\mathfrak F$ be a specified unital $C^*$-algebra and let
$\alpha:\mathfrak F\to\mathfrak F$ be a $*$-automorphism representing one
microscopic update. The algebra must contain the bounded source operators
used below, or an explicitly controlled extension containing them. For
Gauss-dressed charged sources, existence of that extension and compatibility
with the charge sectors are part of the construction. They do not follow
from a formal path label. Starting from a specified state $\omega_{\rm seed}$,
define
\begin{equation}
 \omega_T(A)=\frac1T\sum_{t=0}^{T-1}
       \omega_{\rm seed}(\alpha^{-t}(A)).
 \label{eq:algebraic-cesaro-state}
\end{equation}
The sign agrees with the density average
$T^{-1}\sum_t U^t\rho_{\rm seed}U^{-t}$ when
$\alpha(A)=UAU^\dagger$.

\begin{proposition}[Invariant cluster states]
\label{prop:algebraic-invariant-cluster}
The states $\omega_T$ have weak-$*$ cluster points. Every such point $\omega$
is $\alpha$-invariant, and
\begin{equation}
 |\omega_T(\alpha^n(A))-\omega_T(A)|
 \leq\min\{2,2|n|/T\}\,\|A\|.
 \label{eq:algebraic-stationarity-bound}
\end{equation}
If a symmetry acts by automorphisms commuting with $\alpha$ and leaves the
seed state invariant, it also leaves the averages and their cluster states
invariant.
\end{proposition}
\begin{proof}
The state space of a unital $C^*$-algebra is weak-$*$ compact. Cancellation
of the overlapping terms in the two averages gives the displayed bound;
the universal bound is $2\|A\|$. Taking a convergent subnet proves
invariance. The symmetry statement follows term by term. On a separable
algebra, weak-$*$ compactness of the state space is metrizable, so convergent
subsequences suffice.
\end{proof}

Existence does not imply uniqueness. A specified subnet or invariant mean is
a state-selection choice unless all admissible limits agree on the source
correlation algebra. Unique ergodicity would be a sufficient additional
condition, but is not asserted for QICT. More general ergodic convergence
conditions are studied in Ref.~\cite{Fidaleo2020Ergodic}. The primitive-cost
seed, the stationary average, and a Floquet branch therefore remain distinct
objects. No claim of a gap, purity, positive energy, or a vacuum follows from
Proposition~\ref{prop:algebraic-invariant-cluster} alone.

\subsection{Dynamic GNS representation of the six-source kernel}

Let $O_a\in\mathfrak F$, $a=1,\ldots,6$, be the prescribed dressed sources,
ordered as $L1,L2,L3,R1,R2,R3$. For an invariant state $\omega$, set
\begin{equation}
 K_{ab}(n)=\omega(O_a^*\alpha^n(O_b)),\qquad
 G=K(0),\qquad G\succeq\gamma I_6,\quad\gamma>0.
 \label{eq:algebraic-raw-kernel}
\end{equation}

\begin{theorem}[Stationary source realization]
\label{thm:algebraic-source-realization}
There are a Hilbert space, a unitary $W_\omega$, and an isometry
$V_\omega:\mathbb C^6\to\mathcal H_\omega$ such that
\begin{equation}
 C_n=G^{-1/2}K(n)G^{-1/2}
     =V_\omega^*W_\omega^nV_\omega.
 \label{eq:algebraic-normalized-moments}
\end{equation}
Consequently $C_0=I_6$, $C_{-n}=C_n^*$, and every block Toeplitz matrix
$[C_{j-i}]_{i,j=0}^N$ is positive semidefinite. The source measure
$M(B)=V_\omega^*E_{W_\omega}(B)V_\omega$ is a normalized positive matrix
measure. These statements require no finite-rank Riesz projection of the
microscopic update and no trace-class realization of $\omega$.
\end{theorem}
\begin{proof}
In the GNS representation $(\pi_\omega,\mathcal H_\omega,\Omega_\omega)$,
define on the cyclic dense domain
\[
 W_\omega\pi_\omega(A)\Omega_\omega
     =\pi_\omega(\alpha(A))\Omega_\omega.
\]
Invariance preserves the inner product and the null space. Surjectivity of
$\alpha$ makes the extension unitary. With
$Qe_a=\pi_\omega(O_a)\Omega_\omega$, one has $Q^*Q=G$ and
$Q^*W_\omega^nQ=K(n)$. Taking $V_\omega=QG^{-1/2}$ proves the realization.
Toeplitz positivity follows from
\[
 \sum_{i,j=0}^N v_i^*C_{j-i}v_j
 =\left\|\sum_{j=0}^N W_\omega^jV_\omega v_j\right\|^2\geq0.
\]
The spectral theorem supplies the matrix measure.
\end{proof}

This is a dynamical extension of the static GNS controls in the archive.
The distinction matters: the implementing unitary $W_\omega$ acts in the
representation of the selected stationary functional. It need not have the
spectral type of the original microscopic unitary acting on the Haar carrier.
Its finite source residues do not contradict the Haar finite-rank Riesz
obstruction.

When the finite regulator admits a stationary density $\rho$ on the neutral
sector, the construction agrees with the two-branch prescription. Take the
linking algebra on $\mathcal H_0\oplus\mathcal H_f$,
$U=U_0\oplus U_f$, the state $\operatorname{Tr}(\operatorname{diag}(\rho,0)\,
\cdot)$, and sources in the neutral-to-charged corner. Then
\begin{align}
 K_{ab}(n)
 &=\operatorname{Tr}(\rho O_a^\dagger U_f^nO_bU_0^{-n})\\
 &=\left\langle O_a\sqrt\rho,
           U_f^nO_b\sqrt\rho U_0^{-n}\right\rangle_{\rm HS}.
 \label{eq:algebraic-finite-equivalence}
\end{align}
The second identity uses $[\rho,U_0]=0$. Thus the finite implementation uses
$\mathcal W(X)=U_fXU_0^\dagger$ exactly as before. Replacing it with powers
of $U_0^\dagger U_f$ remains unjustified. If $G$ is singular, its null
directions must be diagnosed or quotiented explicitly; an arbitrary
six-dimensional orthonormal replacement would change the observable.

\subsection{Convergence of observables without trace-norm convergence of states}

For compatible source labels, let $K_R(n)$ be the unnormalized stationary
kernels of a regulator sequence. The condition needed for their spectral
limit is convergence of these functionals on the tested source products.
Global trace-norm convergence of density matrices is a stronger condition
and may be unavailable.

\begin{proposition}[Kernel stability under Gram normalization]
\label{prop:algebraic-kernel-error}
Suppose $G=K(0)$ and $\widetilde G=\widetilde K(0)$ both satisfy
$G,\widetilde G\succeq\gamma I$, and suppose
\[
 \|\widetilde K(n)-K(n)\|\leq\epsilon_n,\quad
 \|\widetilde G-G\|\leq\epsilon_0,\quad \|K(n)\|\leq B_n.
\]
Then the normalized kernels obey
\begin{equation}
 \|\widetilde C_n-C_n\|
 \leq\frac{\epsilon_n}{\gamma}
       +\frac{B_n\epsilon_0}{\gamma^2}.
 \label{eq:algebraic-kernel-error}
\end{equation}
For stationary source kernels one may take $B_n=\|G\|$. At $n=0$,
normalization gives zero difference exactly.
\end{proposition}
\begin{proof}
The resolvent integral for the inverse square root gives
\[
 G^{-1/2}=\frac1\pi\int_0^\infty
 t^{-1/2}(G+tI)^{-1}\,dt,
\qquad
 \|\widetilde G^{-1/2}-G^{-1/2}\|
 \leq\frac{\epsilon_0}{2\gamma^{3/2}}.
\]
Insert the difference of the raw kernels between two factors
$\widetilde G^{-1/2}$, then expand the change of the two normalization
factors around $K(n)$. Since both inverse square roots have norm at most
$\gamma^{-1/2}$, the triangle inequality proves the bound. The realization
$K(n)=Q^*W_\omega^nQ$ gives $\|K(n)\|\leq\|Q\|^2=\|G\|$.
\end{proof}

In particular, suppose $G_R\succeq\gamma I$ uniformly and, for every fixed
integer $n$,
\begin{equation}
 \|K_{R+1}(n)-K_R(n)\|\leq b_{R,n},\qquad
 \sum_R b_{R,n}<\infty.
 \label{eq:algebraic-kernel-cauchy}
\end{equation}
Then $K_R(n)$ has a unique limit along that sequence, with error at most
$\sum_{s\geq R}b_{s,n}$. The proposition controls the corresponding
normalized moments. Toeplitz positivity passes to the limit, and convergence
of every Fourier coefficient gives weak convergence of the normalized
matrix measures on the unit circle. To see the last assertion, first test
trigonometric polynomials and then uniformly approximate a continuous
function; positivity and total mass $I_6$ bound the approximation error.
Finitely many volumes or lags do not prove the summable tails in
Eq.~\eqref{eq:algebraic-kernel-cauchy}. Independence of another admissible
state-selection or regulator sequence requires a comparison on the same
source products.

\subsection{An exact counterexample to the trace-density requirement}

On $\ell^2(\mathbb Z)$, let $S|j\rangle=|j+1\rangle$ and
$D_a|j\rangle=e^{2\pi iaj/7}|j\rangle$ for $a=1,\ldots,6$. Let
$\mathfrak F=C^*(I,S,D_1,\mathcal K(\ell^2(\mathbb Z)))$, where
$\mathcal K$ denotes the compact operators, and let $\alpha=\operatorname{Ad}S$.
The shift has purely continuous spectrum. Its averages from $|0\rangle$ are
\begin{equation}
 \rho_T=\frac1T\sum_{j=0}^{T-1}|j\rangle\langle j|,
 \qquad \|\rho_T-\rho_{2T}\|_1=1.
 \label{eq:algebraic-shift-trace-failure}
\end{equation}
Nevertheless their functionals converge on $\mathfrak F$. On the dense
polynomials in $S,D_1$ and compact operators, the limit is determined by
\[
 \omega(S^mD_1^k)=\begin{cases}1,&m=0,\ k=0\pmod7,\\0,&\text{otherwise},\end{cases}
 \qquad \omega(A)=0\quad(A\in\mathcal K).
\]
Uniform boundedness of states extends convergence to the norm closure.
The limit is stationary and cannot be represented by a trace-class density
on the original $\ell^2(\mathbb Z)$, since it vanishes on every rank-one
projection but has $\omega(I)=1$.

With $O_a=D_a$, one has $\alpha^n(D_b)=e^{-2\pi ibn/7}D_b$ and hence
\begin{equation}
 G=I_6,\qquad
 C_n=\operatorname{diag}_{b=1}^6(e^{-2\pi ibn/7}).
 \label{eq:algebraic-shift-six-moments}
\end{equation}
These six source moments are already exact for $T$ divisible by seven.
Their measure has six atoms, although $S$ has no eigenvectors. The atoms
are eigenvalues of the GNS implementing unitary on the source span, not
finite-rank spectral projections of $S$. This example demonstrates the
mathematical mechanism only. The number seven is an illustrative period;
it is not identified with the QICT receiver, charged particles, or a mass
hierarchy.

\subsection{Coherent C10 dressing and its error budget}

The algebraic state construction retains the prescribed dressing. Let a
finite acyclic directed graph have a unique sink $b$ and nonnegative edge
weights with $\sum_w p_{vw}=1$. For unitary link transports, or more generally
contractions $U_{vw}$ on a common carrier, the ordered dressing obeys
\begin{equation}
 D_b=I,\qquad D_v=\sum_w p_{vw}U_{vw}D_w.
 \label{eq:algebraic-dag-recursion}
\end{equation}
Partitioning paths by their first edge proves that this is precisely
$\sum_{\gamma:v\to b}(\prod_{e\in\gamma}p_e)U_\gamma$. Reverse topological
order proves uniqueness and $\|D_v\|\leq1$. Gauge transformations give
$D_v\mapsto g_vD_vg_b^\dagger$. Application to a vector or source block
requires one link application per edge, while the cost of each application
still depends on the physical Hilbert or tensor dimensions.

If $\|\widetilde U_e-U_e\|\leq\epsilon_e$ and both exact and approximate
links are contractions, the same recursion and the triangle inequality give
\begin{equation}
 \|\widetilde D_v-D_v\|\leq\sum_e F_v(e)\epsilon_e,
 \label{eq:algebraic-dag-error}
\end{equation}
where $F_v(e)$ is the expected occupation of edge $e$ in the prescribed path
measure. For the C10 positive-flow decomposition these are its conserved
edge flows. Explicitly, the error recursion is
$a_b=0$, $a_v=\sum_w p_{vw}(\epsilon_{vw}+a_w)$; its expansion is the
expected accumulated path error. If
$\|\widetilde U_e\|\leq s_e$ instead, use
$a_v=\sum_w p_{vw}(\epsilon_{vw}+s_{vw}a_w)$.
Probability uncertainties, arithmetic rounding, and approximate state
functionals require separate bounds.

Coherent terms between distinct paths occur in $D\rho D^\dagger$.
Replacing this expression by a Kraus mixture with operators
$\sqrt{P(\gamma)}U_\gamma$ deletes those terms and changes the source.
Likewise, intertwiner multiplicities remain coherent coordinates unless an
independent superselection argument applies. The recursion supplies the
dressing once the link actions are available; it does not generate the
missing interacting state or global matter--gauge update.

\paragraph{Archived geometric-flow check.}
The $L=24$, $n=9$ precontraction supplies 3925 interior vertices and 17576
oriented edges, with the boundary identified as one sink. Applying the
recursion to its recorded probabilities gives a mean path length
$13.217861267754989$ in binary64 arithmetic. The resulting edge occupations
agree with the archived flows to $1.79\times10^{-15}$ in maximum absolute
error; the conservation residual is below $4.21\times10^{-15}$.
For a uniform certified link error $\epsilon$ and contraction-valued links,
Eq.~\eqref{eq:algebraic-dag-error} therefore has the geometric factor
$\sum_e F(e)$, numerically about $13.218$. These decimals are evaluations,
not outward-rounded enclosures. Identity link callbacks were used solely
to check graph transport and its adjoint. They do not provide the
interacting quantum links or a physical C10 kernel. The input checksum,
code and residuals are supplied with the computational archive.

\subsection{What is resolved, and what the physical limit requires}

The source observable now has a consistent stationary formulation in domains
where finite-rank Floquet selection is obstructed. Its finite normal-state
restriction agrees with the existing relative moments. Observable convergence
can be tested through Eq.~\eqref{eq:algebraic-kernel-cauchy}, rather than an
unavailable global trace-density limit. The positive-moment, resolvent, and
finite-resolution spectral estimates therefore remain applicable in the
chosen stationary representation.

Physical completion still requires the full ordered interacting update,
Gauss-compatible dynamical sources, and a selected limiting functional on
their algebra. A stationary state need not be a vacuum. A vacuum
interpretation additionally requires a compatible continuous-time physical
limit with a positive-energy generator and the stated reference-energy
condition. Particle identification requires the appropriate source overlap,
momentum dependence, branch and clock calibration, and threshold or
resonance analysis. The unit-circle limit alone does not justify unwrapping
phases into physical energies as the regulator tick tends to zero.

In particular, a uniform gap above the physical vacuum is not a universal
requirement for a theory containing massless excitations. An auxiliary
primitive-cost gap, a finite Riesz contour, and the physical excitation gap
are different hypotheses. Charged infrared thresholds and gravitational
massless spin-two response retain their separate conditions. Nothing in the
GNS construction determines the CKM or PMNS matrices, a dark-particle
abundance, or the Einstein-response coefficient without their interacting
correlations. The existing physical validation requirements are retained;
this appendix supplies an alternative mathematical route that must itself
be instantiated and checked.

The accompanying tests verify the finite normal-state equivalence, the exact
shift control, kernel normalization, the perturbation inequality, and the
coherent path recursion. Floating-point checks exercise those implementations;
they are not numerical QICT particle predictions or interval certificates.
